% Part: first-order-logic
% Chapter: sequent-calculus
% Section: rules-and-proofs

\documentclass[../../../include/open-logic-section]{subfiles}

\begin{document}

\olfileid{mvl}{seq}{rul}

\olsection{规则与 \usetoken{P}{derivation}}

在以下讨论中，令 $\Gamma, \Delta, \Pi, \Lambda$ 表示有穷的语句序列。

\begin{defn}[相继式]
一个\emph{$n$ 边相继式}是形如
\[
\Gamma_1 \nSequent \dots \nSequent \Gamma_n
\]
的表达式，其中每个 $\Gamma_1$ 是语言~$\Lang L$ 中语句的一个有穷（可能为空）序列。
\end{defn}

\begin{defn}[初始相继式]
对于语言中的任意语句 $!A$，形如 $!A \nSequent \dots \nSequent !A$ 的 $n$ 边相继式称为\emph{$n$ 边初始相继式}。

如果语言包含一个 0 元联结词~$\star$，即一个命题常项，那么我们也把形如 $\dots \nSequent \star \nSequent \dots$ 的相继式作为初始相继式，其中 $\star$ 出现在与其真值~$\tf{\star} \in V$ 相关联的位置上，而其他位置为空。
\end{defn}

对于 $n$ 值逻辑~$\Log{L}$ 的每个联结词，该联结词在~$\Log{L}$ 中可取的每个真值都对应一条逻辑规则。在 $\Log{L}$ 的 $n$ 边相继式演算中，推导是由相继式构成的树，其最顶层的相继式是初始相继式；并且，若一个相继式位于一个或多个其他相继式之下，则它必须是由 $\Log{L}$ 联结词的某条推理规则正确生成的。

\begin{defn}[定理]
如果存在一个 $n$ 边相继式的推导，使得该相继式在每个对应于 $\Log{L}$ 的指定真值的位置上都包含 $!A$，则称 $n$ 值逻辑~$\Log{L}$ 中的语句~$!A$ 为 $\Log{L}$ 的一个\emph{定理}。
若 $!A$ 是定理，则记作 $\Proves[\Log{L}] !A$；否则记作 $\Proves/[\Log{L}] !A$。
\end{defn}

\begin{defn}[可推导性]
在 $n$ 值逻辑~$\Log{L}$ 中，语句~$!A$ \emph{可从}语句集合~$\Gamma$ \emph{推导}，记作 $\Gamma \Proves[\Log{L}] !A$，当且仅当存在 $\Gamma$ 的一个有穷子集~$\Gamma_0 \subseteq \Gamma$ 以及 $\Gamma_0$ 中语句的一个序列 $\Gamma_0'$，使得如下相继式存在推导：
\[ \Lambda_1 \nSequent \dots \nSequent \Lambda_n \]
其中 $\Lambda_i$ 在位置 $i$ 对应于指定真值时为 $!A$，否则为 $\Gamma_0'$。如果 $!A$ 不能从 $\Gamma$ 推导，则记作 $\Gamma \Proves/ !A$。
\end{defn}

例如，三值 Łukasiewicz（卢卡西维茨）逻辑有一个三边相继式演算。在三边相继式 $\Gamma \nSequent \Pi \nSequent \Delta$ 中，$\Gamma$ 对应于~$\False$，$\Delta$ 对应于~$\True$，而 $\Pi$ 对应于~$\Undef$。公理是 $!A \nSequent !A \nSequent !A$。由于只有 $\True$ 是指定真值，因此 $\Gamma \Proves[\LogLuk[3]] !A$ 当且仅当相继式 $\Gamma \nSequent \Gamma \nSequent !A$ 存在推导。（如果 $\Undef$ 也是指定真值，我们就需要 $\Gamma \nSequent !A \nSequent !A$ 的推导。）

\end{document}